\documentclass{controle}
\usepackage{main}

\title{Activité : Jeu du Juniper Green}
\author{Seconde}
\date{Chapitre 7 Partie 1}

\begin{document}
\maketitle

Le jeu du Juniper Green est un jeu entre deux joueurs nécessitant un tableau de tous les nombres de $1$ à $50$.

\begin{center}
\begin{tblr}{hlines,vlines,columns={wd=0.8cm,c},rows={ht=1cm,m}}
1&2&3&4&5&6&7&8&9&10\\
11&12&13&14&15&16&17&18&19&20\\
21&22&23&24&25&26&27&28&29&30\\
31&32&33&34&35&36&37&38&39&40\\
41&42&43&44&45&46&47&48&49&50\\
\end{tblr}
\end{center}
Les règles sont les suivantes :
\begin{enumerate}
\item Le premier joueur choisit un nombre pair et le coche;
\item À tour de rôle, le joueur dont c'est le tour choisit un multiple ou un diviseur du nombre venant d'être coché. \textbf{Il ne peut pas choisir un nombre déjà coché durant la partie.}
\item Le joueur dont c'est le tour mais qui ne peut plus jouer a perdu.
\end{enumerate}

\begin{questions}
\question 
Jouer quelques parties du jeu avec votre voisin; puis répondre aux questions suivantes :
\begin{parts}
\part Sans tenir compte de la règle 1.; si on suppose que le joueur $1$ coche la case $1$, quelle case le joueur suivant peut-il cocher ?
\part Sans tenir compte de l règle 1.; quelle case le joueur $1$ doit-il cocher en premier pour que le joueur suivant aie le moins de choix possibles ? 
\part Décrire une stratégie gagnante pour le premier joueur si on ne tient pas compte de la règle 1.
\part Quelle est l'utilité de la règle 1. ? Y a-t-il toujours une stratégie gagnante pour le premier joueur ?
\end{parts}
\end{questions}
\end{document}